\clearpage
\section{Requirements}

Sebanyak 24 kebutuhan sistem dikelompokkan berdasarkan fungsi, kinerja, keselamatan, antarmuka, keandalan, dan batasan desain. Setiap blok di bawah memisahkan pernyataan kebutuhan, sumber, alasan, metode verifikasi, dan penelusuran model.

Pada baris sumber, EVID-* menunjuk evidence dan ASM-* menunjuk asumsi desain. Angka baseline tetap merupakan asumsi. Metode Test, Analysis, dan Inspection berarti uji, analisis, dan inspeksi; pencantuman metode tidak menyatakan bahwa pemeriksaan sudah lulus.

Penelusuran dibaca sebagai use case (UC) $\rightarrow$ aksi (ACT) $\rightarrow$ blok (BLK) $\rightarrow$ skenario verifikasi (VER). ID yang sama digunakan pada diagram dan matriks traceability.

\Needspace{16\baselineskip}
\subsection{Fungsional (F)}

\requirementcard
{REQ-F-001 --- Klasifikasi}
{DensityAnalysis harus memetakan count valid ke LOW (0..5), MEDIUM (6..15), atau HIGH ($\geq$16).}
{EVID-TRF-006, ASM-003}
{Proksi count tanpa kapasitas lokasi.}
{Test / VER-01}
{UC-01 $\rightarrow$ ACT-02 $\rightarrow$ BLK-03 $\rightarrow$ VER-01.}

\requirementcard
{REQ-F-002 --- Eligibility}
{TrafficController harus mengeluarkan pendekat dari eligible setelah grant green hingga round berikutnya.}
{ASM-004}
{Satu pelayanan per pendekat per round.}
{Analysis / VER-02}
{UC-02 $\rightarrow$ ACT-08 $\rightarrow$ BLK-04 $\rightarrow$ VER-02.}

\requirementcard
{REQ-F-003 --- Prioritas}
{TrafficController harus memilih eligible menurut kategori menurun, waktu tunggu menurun, lalu jarak ring dari pointer bila deadline nominal masih layak.}
{EVID-TRF-001, ASM-004}
{Adaptif dengan tie deterministik; bukan formula jurnal.}
{Test / VER-03}
{UC-02 $\rightarrow$ ACT-03 $\rightarrow$ BLK-04 $\rightarrow$ VER-03.}

\requirementcard
{REQ-F-004 --- Grant green}
{TrafficController harus mengunci grant 10/20/30 s untuk LOW/MEDIUM/HIGH hingga green berakhir pada mode Normal.}
{EVID-TRF-001, ASM-005, ASM-006}
{Kategori lebih tinggi mendapat grant lebih lama tanpa perubahan di tengah fase.}
{Test / VER-04}
{UC-02 $\rightarrow$ ACT-04 $\rightarrow$ BLK-04 $\rightarrow$ VER-04.}

\requirementcard
{REQ-F-005 --- Timer tunggu}
{TrafficController harus mengukur wait sejak green terakhir berakhir, nol selama green dan sejak startup untuk pendekat belum dilayani.}
{ASM-004}
{Definisi service wait tidak sama dengan delay individu.}
{Test / VER-05}
{UC-02 $\rightarrow$ ACT-03 $\rightarrow$ BLK-04 $\rightarrow$ VER-05.}

\Needspace{16\baselineskip}
\subsection{Kinerja dan batas waktu (P)}

\requirementcard
{REQ-P-001 --- Observasi}
{VehicleDetection harus menerima frame nominal lengkap pada cadence 1 s di model.}
{ASM-002, ASM-006}
{Jendela observasi dan sampling eksplisit.}
{Test / VER-06}
{UC-01 $\rightarrow$ ACT-01 $\rightarrow$ BLK-02 $\rightarrow$ VER-06.}

\requirementcard
{REQ-P-002 --- Boundary keputusan}
{TrafficController harus menerbitkan grant berikutnya hanya pada tick 1 s di akhir all-red.}
{ASM-006}
{Semua durasi pada grid diskrit yang sama.}
{Analysis / VER-07}
{UC-02 $\rightarrow$ ACT-03 $\rightarrow$ BLK-04 $\rightarrow$ VER-07.}

\requirementcard
{REQ-P-003 --- Minimum green}
{SignalController harus mempertahankan green yang diterima sekurang-kurangnya 10 s kecuali masuk SafetyHold.}
{ASM-005, ASM-008}
{Batas nominal dan pengecualian keselamatan eksplisit.}
{Test / VER-08}
{UC-03 $\rightarrow$ ACT-05 $\rightarrow$ BLK-05 $\rightarrow$ VER-08.}

\requirementcard
{REQ-P-004 --- Maximum green}
{SignalController harus mengakhiri setiap grant green paling lambat 30 s setelah grant dimulai.}
{ASM-005}
{Membatasi layanan tunggal.}
{Analysis / VER-08}
{UC-03 $\rightarrow$ ACT-05 $\rightarrow$ BLK-05 $\rightarrow$ VER-08.}

\requirementcard
{REQ-P-005 --- Maximum waiting}
{TrafficController harus menjadwalkan setiap pendekat kembali green dalam 210 s sejak green sebelumnya berakhir pada operasi round baseline tanpa SafetyHold/restart.}
{ASM-004, ASM-005, ASM-008}
{Bound analitis nominal 208 s.}
{Analysis / VER-09}
{UC-02 $\rightarrow$ ACT-03 $\rightarrow$ BLK-04 $\rightarrow$ VER-09.}

\Needspace{16\baselineskip}
\subsection{Keselamatan (S)}

\requirementcard
{REQ-S-001 --- Eksklusivitas}
{SignalController harus menghasilkan paling banyak satu pendekat green pada setiap state model.}
{EVID-TRF-002, ASM-001, ASM-008}
{Matriks konflik konservatif.}
{Analysis / VER-10}
{UC-03 $\rightarrow$ ACT-05 $\rightarrow$ BLK-05 $\rightarrow$ VER-10.}

\requirementcard
{REQ-S-002 --- Yellow}
{SignalController harus menjalankan yellow 3 s pada pendekat aktif sebelum clearance all-red setelah grant normal berakhir.}
{ASM-008}
{Pemisahan transisi eksplisit; bukan waktu lokasi terkalibrasi.}
{Test / VER-11}
{UC-03 $\rightarrow$ ACT-06 $\rightarrow$ BLK-05 $\rightarrow$ VER-11.}

\requirementcard
{REQ-S-003 --- All-red}
{SignalController harus mempertahankan semua pendekat merah selama 1 s pada startup dan setiap clearance sebelum green berikutnya.}
{ASM-008}
{Batas fase aman.}
{Test / VER-11}
{UC-03 $\rightarrow$ ACT-06 $\rightarrow$ BLK-05 $\rightarrow$ VER-11.}

\requirementcard
{REQ-S-004 --- Command rejection}
{SignalController harus masuk SafetyHold all-red bila menerima command di luar domain atau command green yang melompati clearance.}
{ASM-008}
{Tidak mengeksekusi command yang melanggar invariant.}
{Test / VER-12}
{UC-03 $\rightarrow$ ACT-09 $\rightarrow$ BLK-05 $\rightarrow$ VER-12.}

\requirementcard
{REQ-S-005 --- Deadline precedence}
{TrafficController harus mengganti urutan nominal dengan earliest-deadline-first jika proyeksi sisa round nominal melampaui deadline 210 s.}
{ASM-004, ASM-005}
{Fairness mendahului density pada boundary.}
{Test / VER-13}
{UC-02 $\rightarrow$ ACT-03 $\rightarrow$ BLK-04 $\rightarrow$ VER-13.}

\requirementcard
{REQ-S-006 --- Infeasible plan}
{TrafficController harus meminta SafetyHold bila urutan earliest-deadline-first juga tidak memenuhi deadline.}
{ASM-004, ASM-008}
{Jangan menjanjikan bound saat keadaan tidak konsisten.}
{Test / VER-13}
{UC-03 $\rightarrow$ ACT-03 $\rightarrow$ BLK-04 $\rightarrow$ VER-13.}

\Needspace{16\baselineskip}
\subsection{Antarmuka (I)}

\requirementcard
{REQ-I-001 --- Input frame}
{VehicleDetection harus memvalidasi CountsFrame sesuai domain field, empat ID unik dan aturan timestamp pada kontrak operasi.}
{ASM-002, ASM-006}
{Input tervalidasi sebelum analisis.}
{Test / VER-14}
{UC-01 $\rightarrow$ ACT-01 $\rightarrow$ BLK-02 $\rightarrow$ VER-14.}

\requirementcard
{REQ-I-002 --- Output lampu}
{TrafficLights harus menerima LampVector empat indikasi yang masing-masing bernilai Red, Yellow atau Green dari SignalController.}
{EVID-TRF-002, ASM-008}
{Kontrak output logis; actuator ideal.}
{Inspection / VER-15}
{UC-03 $\rightarrow$ ACT-05 $\rightarrow$ BLK-06 $\rightarrow$ VER-15.}

\Needspace{16\baselineskip}
\subsection{Keandalan dan penanganan kegagalan (R)}

\requirementcard
{REQ-R-001 --- Deteksi data fault}
{VehicleDetection harus melatch dataFault ketika frame invalid, valid=false, cache tidak lengkap atau umur cached sample lebih dari 3 s.}
{ASM-006, ASM-007}
{Data rusak tidak diperlakukan sebagai count nol.}
{Test / VER-16}
{UC-04 $\rightarrow$ ACT-01 $\rightarrow$ BLK-02 $\rightarrow$ VER-16.}

\requirementcard
{REQ-R-002 --- Masuk fallback}
{TrafficController harus memakai Fallback pada boundary aman berikutnya setelah dataFault dilatch tanpa mengubah grant aktif.}
{ASM-007}
{Tidak memotong green oleh fault data.}
{Test / VER-17}
{UC-04 $\rightarrow$ ACT-07 $\rightarrow$ BLK-04 $\rightarrow$ VER-17.}

\requirementcard
{REQ-R-003 --- Fallback timing}
{TrafficController harus melayani eligible dalam urutan ring dengan green 10 s pada Fallback.}
{ASM-007}
{Timing default deterministik; round tidak direset.}
{Test / VER-18}
{UC-04 $\rightarrow$ ACT-07 $\rightarrow$ BLK-04 $\rightarrow$ VER-18.}

\requirementcard
{REQ-R-004 --- Pemulihan}
{TrafficController harus kembali Normal hanya pada awal round setelah tiga frame lengkap fresh/valid dengan timestamp baru berturut-turut di tick berbeda.}
{ASM-007}
{Mencegah flapping dan pemulihan dari retransmisi yang sama.}
{Test / VER-19}
{UC-04 $\rightarrow$ ACT-07 $\rightarrow$ BLK-04 $\rightarrow$ VER-19.}

\Needspace{16\baselineskip}
\subsection{Batasan desain (D)}

\requirementcard
{REQ-D-001 --- Scope}
{Model harus membatasi operasi pada empat pendekat lurus tanpa pedestrian, dedicated turn, emergency atau ML.}
{ASM-001, ASM-009}
{Menjaga scope pengguna.}
{Inspection / VER-20}
{UC-02 $\rightarrow$ ACT-03 $\rightarrow$ BLK-01 $\rightarrow$ VER-20.}

\requirementcard
{REQ-D-002 --- SysML}
{Model harus mempertahankan identitas dan semantik subset SysML 1.6 dalam tujuh view dan pemetaan renderer yang terdokumentasi.}
{EVID-LEC-001, EVID-LEC-002, EVID-LEC-003, EVID-LEC-004, ASM-010}
{Semantik tidak ditentukan tool rendering.}
{Inspection / VER-21}
{UC-02 $\rightarrow$ ACT-03 $\rightarrow$ BLK-01 $\rightarrow$ VER-21.}

\clearpage
